\scanpage{22}
\mspage{10}
\margin{$(g,\nu=\operatorname{card}J)$ anabélien}
\uline{NB} Si $\operatorname{card}J=\nu$, alors $\mathcal M_{g,J}$ est isomorphe au schéma modulaire $\mathcal M^!_{g,\nu}$ (via un iso. $[1,\nu]\cap\mathbb Z\simeq J$), de façon compatible avec les autom. des groupes $\unsure{\mathfrak S_J|\mathfrak G_J}$ et $\unsure{\mathfrak S_\nu|\mathfrak G_\nu}$. On sait que $\mathcal M^!_{g,\nu}$ est lisse de type fini sur $\mathbb Z$. Il en est donc de $=$ de $\mathcal M_{g,J}$, donc aussi de $\mathcal M_G$ (pour $G$ semi-stable) et de $\mathcal M_{[G]}=(\mathcal M_G,\Gamma)$.

4. \uline{La théorie de Mumford-Deligne}

Soient $\mathcal S$ une multiplicité schématique, $\unsure{\mathcal X|\mathfrak X}$ une \deleted{courbe} schéma relatif sur $\mathcal S$, propre sur $\mathcal S$, $\underline I$ un sous-schéma fermé de $\mathcal X$. On dit que $(\mathcal X,\underline I)$ est une MD-courbe relative sur $\mathcal S$, si $\mathcal X,\underline I$ sont plats de présentation finie sur $\mathcal S$, et si pour $\forall$ pt géom. de $\mathcal S$, la
\mspage{21}
